\chapter{Background}
\label{ch:background}

This chapter sets out the concepts the rest of the report relies on: the MPS domain, the software licensing categories, state-of-the-practice methodology, the nine quality attributes, commonality analysis, the Analytic Hierarchy Process, and the role of repository-based measurement.

\section{Robotics Software for Motion Planning and Simulation}
\label{sec:robotics-software-background}

MPS software supports the design, testing, evaluation, and deployment of robotic systems. Motion planning software generates feasible robot motions under constraints such as geometry, kinematics, dynamics, collisions, and task requirements~\citep{lavalle2006planning}. Simulation software provides virtual environments for representing robots, sensors, actuators, controllers, and physical interactions before, or alongside, work with physical hardware~\citep{collins2021simulators}.

The MPS domain in this report is defined to span more than a single class of software, so that the study includes enough packages to support a meaningful comparison. It covers full simulation environments, physics and dynamics engines, motion planning libraries, manipulation frameworks, and middleware. These classes of software are not interchangeable, but they are unified because they appear together in the same workflow: a simulator may call a physics engine for contact and dynamics, a planning stack may sit on top of a collision-checking library and communicate through middleware, and middleware may call planning, control, perception, and simulation. Figure~\ref{fig:mps-software-stack} sketches these relationships among the four software roles; the report is explicit throughout about which role each package plays.

\begin{figure}[htbp]
\centering
% Figure palette: the two structural rails (middleware, physics) share one
% neutral slate; the two subject roles (planning, simulation) each carry one
% accent. Three colours total, each with a matching deeper border.
\definecolor{mpsSlateFill}{HTML}{EDEEF1}
\definecolor{mpsSlateLine}{HTML}{5E6675}
\definecolor{mpsSlateText}{HTML}{394150}
\definecolor{mpsTealFill}{HTML}{E0F1EC}
\definecolor{mpsTealLine}{HTML}{0F8E78}
\definecolor{mpsTealText}{HTML}{0A5446}
\definecolor{mpsBlueFill}{HTML}{E6EEFA}
\definecolor{mpsBlueLine}{HTML}{2C6FCE}
\definecolor{mpsBlueText}{HTML}{1A4A8F}
\definecolor{mpsArrow}{HTML}{7A828C}
\begin{tikzpicture}[
  font=\small,
  comp/.style={draw, rounded corners, align=center, inner sep=5pt, minimum height=1.15cm, line width=0.8pt},
  mw/.style={comp, fill=mpsSlateFill, draw=mpsSlateLine, text=mpsSlateText, text width=9cm},
  plan/.style={comp, fill=mpsTealFill, draw=mpsTealLine, text=mpsTealText, text width=4cm},
  sim/.style={comp, fill=mpsBlueFill, draw=mpsBlueLine, text=mpsBlueText, text width=4cm},
  phys/.style={comp, fill=mpsSlateFill, draw=mpsSlateLine, text=mpsSlateText, text width=9cm},
  bus/.style={{Latex[length=2mm]}-{Latex[length=2mm]}, thick, draw=mpsArrow},
  uses/.style={-{Latex[length=2.6mm]}, thick, draw=mpsArrow},
]
\node[mw] (mw) at (2.5,0) {\textbf{Middleware} \footnotesize(ROS~2, YARP)\\[1pt]\footnotesize communication, hardware abstraction, system integration};
\node[plan] (plan) at (0,-2.25) {\textbf{Motion planning \&}\\\textbf{manipulation}\\[1pt]\footnotesize OMPL, MoveIt!};
\node[sim] (sim) at (5,-2.25) {\textbf{Simulation}\\\textbf{environments}\\[1pt]\footnotesize Gazebo, Webots, CoppeliaSim, CARLA};
\node[phys] (phys) at (2.5,-5.15) {\textbf{Physics \& dynamics engines}\\[1pt]\footnotesize Bullet, MuJoCo, ODE, DART};
\draw[bus] (plan.north) -- (mw.south -| plan.north);
\draw[bus] (sim.north) -- (mw.south -| sim.north);
\draw[uses] (plan.south) -- node[left=3pt, pos=0.48, font=\footnotesize, align=center, fill=white, inner sep=1pt] {collision /\\dynamics} (phys.north -| plan.south);
\draw[uses] (sim.south) -- node[right=3pt, pos=0.48, font=\footnotesize, align=center, fill=white, inner sep=1pt] {contact /\\dynamics} (phys.north -| sim.south);
\end{tikzpicture}
\caption[Representative MPS software roles]{Representative MPS software roles. Box contents are example packages from the selected set.}
\label{fig:mps-software-stack}
\end{figure}

\section{Software Categories}
\label{sec:software-categories}

Assessing software packages requires understanding the licensing model under which each package is distributed. In particular, the availability of source code determines whether repository-based measurement is possible. Following the terminology used in prior studies of the state of the practice~\citep{Smith2024MI,Smith2024LBM}, this report distinguishes three common software categories:

\begin{description}
    \item[Open Source Software (OSS):] For OSS, the source code is openly accessible. Users have the right to study, change, and distribute it under a license granted by the copyright holder. For many OSS projects, the development process relies on collaboration among contributors worldwide. Accessible source code exposes the underlying logic of software functions, the development practices used to produce them, and the flaws and potential risks in the final product.
    \item[Freeware:] Freeware is software that can be used free of charge. Unlike OSS, the authors of freeware do not permit access to or modification of the source code. To many end users, the distinction between freeware and OSS may not matter.
    \item[Commercial Software:] Commercial software is developed and sold by companies as a product. Typically, commercial software requires payment to access all features and does not include access to the source code. Some commercial software is available free of charge, but its source code is still not accessible.
\end{description}

This report assesses only software that fits under the Open Source Software category. This restriction is a methodological necessity, not a judgement about the quality or importance of freeware or commercial tools.

\section{State-of-the-Practice Studies}
\label{sec:state-of-practice-background}

A study of the state of the practice examines existing tools, artifacts, processes, and evidence in a chosen domain. The goal is not to build new software but to learn from what already exists, by collecting and organizing public evidence in a systematic way. Claims about development activity, documentation, issue management, installation behaviour, or maintainability must be backed by observable project artifacts. When evidence is missing or ambiguous, the gap is recorded rather than filled with assumptions.

The methodology used in this report is based on ``Methodology for Assessing the State of the Practice for Domain X'' by \citet{Smith2021}. It is generic enough to apply to any research software domain and is designed to be completed by a small team within a bounded time budget (Section~\ref{subsec:measurement-execution}). The prescribed process runs from domain selection through candidate filtering, template-based measurement, developer interviews, and AHP-based ranking to answers to a set of research questions; Chapter~\ref{ch:methodology} enumerates the steps as applied to MPS.

The methodology builds on prior assessments of the state of the practice conducted across several research software domains: Geographic Information Systems~\citep{smith2018state}, Mesh Generators~\citep{smith2016state}, Seismology software~\citep{Smith2018Seismology}, and Statistical software for psychology~\citep{smith2018statistical}. More recently, the methodology was refined and applied to Medical Imaging (MI) software~\citep{Smith2024MI} and Lattice Boltzmann Method (LBM) software~\citep{Smith2024LBM,Michalski2021}. These studies are the primary structural and stylistic references for this report.

Practitioner-facing research software checklists provide complementary guidance on themes such as source control, licensing, documentation, testing, automation, distribution, and governance~\citep{Acton2026Checklists}. These themes overlap with several Domain~X quality categories, but this report uses the Domain~X template as its measurement instrument.

\section{Software Quality Attributes}
\label{sec:software-quality-attributes}

This report treats software quality operationally: a package exhibits quality to the extent that observable attributes support dependable installation, use, modification, reuse, and understanding. Following common practice in software engineering, quality is decomposed into distinct attributes, often called ``ilities''~\citep{Smith2021}. This report assesses nine quality attributes drawn from the measurement template in Appendix~B of \citet{Smith2021}; this template is reproduced, as adapted for this study, in Appendix~\ref{appendix:measurement-template} of this report. Several of the qualities use the word ``surface'' to indicate that the assessment is based on shallow but systematic checks performed within a limited time budget, not on exhaustive experimental evaluation~\citep{Smith2024MI}.

\begin{description}
    \item[Installability:] The effort required for the installation and uninstallation of software in a specified environment. Measured through the presence and quality of installation instructions, the number of steps required, the availability of automation, the handling of dependencies, and the presence of validation procedures.
    \item[Correctness and Verifiability:] A program is correct if it matches its specification, which may be stated explicitly or implicitly. The related quality of verifiability is the ease with which the software components or the integrated product can be checked to demonstrate correctness. Measured through the presence of requirements or theory documentation, techniques used to build confidence in correctness (such as automated testing, assertions, or continuous integration), and the availability and quality of tutorials with expected outputs.
    \item[Surface Reliability:] The probability of failure-free operation of a computer program in a specified environment for a specified time. Surface reliability is assessed through whether the software breaks during installation or initial tutorial use, whether descriptive error messages are displayed, and whether recovery is possible.
    \item[Surface Robustness:] Software possesses robustness if it behaves reasonably when it encounters circumstances not anticipated in the requirements specification, and when users violate the assumptions in its requirements specification. Surface robustness is assessed through whether the software handles unexpected or malformed input gracefully, including edge cases such as modified line endings in text input files.
    \item[Surface Usability:] The extent to which a product can be used by specified users to achieve specified goals with effectiveness, efficiency, and satisfaction in a specified context of use. Surface usability is assessed through the presence of getting-started tutorials, user manuals, documented user characteristics, and the availability and variety of user support channels.
    \item[Maintainability:] The effort with which a software system or component can be modified to correct faults, improve performance or other attributes, or satisfy new requirements. Assessed through version information, contribution and review guidance, available non-code artifacts, issue tracking practices, the percentage of identified issues that are closed, the percentage of code that consists of comments, and the version control system in use.
    \item[Reusability:] The extent to which a software component can be used with or without adaptation in a problem solution other than the one for which it was originally developed. Assessed through the number of code files and the presence and quality of API documentation.
    \item[Surface Understandability:] The capability of the software product to enable the user to understand whether the software is suitable, and how it can be used for particular tasks and conditions of use. Surface understandability is assessed through a review of randomly selected source files, examining indentation consistency, coding standards, identifier quality, use of named constants, comment clarity, algorithm references, parameter ordering, and modularization.
    \item[Visibility and Transparency:] The extent to which all steps of a software development process and the current status of that process are conveyed clearly to external observers. Assessed through whether the development process is defined and documented, whether the development environment is documented, and whether release notes are published.
\end{description}

These nine qualities organize public evidence into comparable categories. They are not complete evaluations of each package. Scores may vary with the use case, platform, user background, or measurement date. Later chapters therefore report both the measurements and their limitations.

\section{Commonality Analysis}
\label{sec:commonality-analysis-background}

A domain analysis consists of systematically identifying and documenting the commonalities, variabilities, and parameters of variation of a software family~\citep{Weiss1998}. A software family is defined as ``a set of programs whose common properties are so extensive that it is advantageous to study the common properties of the programs before analyzing individual members''~\citep{Parnas1976}. \citet{Weiss1998} defines commonality analysis as an approach to defining a family by identifying three kinds of information:

\begin{description}
    \item[Commonalities:] Goals, theories, models, definitions, and assumptions that are shared across family members. For example, all robotics simulators share the goal of representing robot kinematics and dynamics in a virtual environment.
    \item[Variabilities:] Goals, theories, models, definitions, and assumptions that differ between family members. For example, simulators may differ in the physics engines they use, the sensor models they support, or the robot platforms they can represent.
    \item[Parameters of Variation:] The possible values for each variability, along with their binding time. Binding time records when a particular variation is fixed: at specification time (when the software is designed), at build time (when the software is compiled), at configuration time (when the software is set up for use), or at run time (when the software is executing).
\end{description}

Commonality analysis was added to the Domain~X methodology after early applications revealed that software packages within a domain can vary considerably, even when they appear to belong to the same family~\citep{Smith2024LBM}. Without it, comparisons risk treating fundamentally different types of software as if they were interchangeable. The analysis records where the packages do similar things and where their roles diverge, so a low score on, for example, installability means different things for a full simulator and a focused dynamics library.

Chapter~\ref{ch:selected-software-commonality} presents this commonality analysis for the selected packages.

\section{Analytic Hierarchy Process}
\label{sec:ahp-background}

The Domain~X methodology ranks the 28 software packages with the Analytic Hierarchy Process (AHP), a pairwise-comparison technique introduced by \citet{Saaty1980}. Each package has one impression score, from 1 to 10, for each of the nine qualities. AHP combines these scores into a single overall priority per package, by which the packages are ordered, following the approach of \citet{smith2016state}.

AHP is applied to one quality at a time. For a given quality, let the package scores be $s_1,\dots,s_n$. Each ordered pair of packages $(i,j)$ is given a comparison value on Saaty's 1--9 scale, read directly from the two scores:
\begin{equation}
a_{ij} =
\begin{cases}
\min\{9,\; s_i - s_j + 1\}, & s_i \ge s_j,\\[4pt]
1 \big/ \min\{9,\; s_j - s_i + 1\}, & s_i < s_j.
\end{cases}
\label{eq:ahp-pairwise}
\end{equation}
Here $a_{ij} = 1$ means the two packages are judged equal on the quality, and $a_{ij} = 9$ is the maximum of the scale, used when package~$i$'s score far exceeds package~$j$'s. The values form an $n \times n$ reciprocal matrix $A$ (with $n = 28$ packages), with $a_{ji} = 1/a_{ij}$. The priority of each package on the quality follows the standard AHP averaging method~\citep{Saaty1980}: each column of $A$ is normalized to sum to one, and the entries of each row are averaged,
\begin{equation}
w_i = \frac{1}{n} \sum_{j=1}^{n} \frac{a_{ij}}{\sum_{l=1}^{n} a_{lj}}.
\label{eq:ahp-priority}
\end{equation}
The priorities $w_1,\dots,w_n$ are non-negative and sum to one, and $w_i$ is the priority of package~$i$ on that quality; sorting the packages by $w_i$ gives their ranking on that quality. Repeating the procedure for all nine qualities gives nine priority vectors $\mathbf{w}^{(1)},\dots,\mathbf{w}^{(9)}$, and hence nine per-quality rankings.

Combining the nine per-quality rankings into a single overall ranking requires weights $c_1,\dots,c_9$ that express the importance of each quality; these weights are obtained by applying the same AHP pairwise-comparison procedure to the qualities themselves. A package's overall priority is the weighted sum of its per-quality priorities:
\begin{equation}
p_i = \sum_{k=1}^{9} c_k\, w_i^{(k)},
\qquad \sum_{k=1}^{9} c_k = 1,
\label{eq:ahp-aggregate}
\end{equation}

The packages are ranked by $p_i$.

AHP is an analysis aid for cross-quality comparison under a stated weighting, not a universal recommendation of any particular package. Selecting a software package for a specific application would involve finding the weights appropriate for that application. The weighting adopted in this report, and the resulting ranking, are stated in Section~\ref{subsec:ahp-ranking} and Chapter~\ref{ch:measurement-results}.

\section{Open Source and Repository-Based Measurement}
\label{sec:open-source-repository-measurement}

The measurement process in this report depends on evidence from public artifacts. Public repositories can reveal commit history, release activity, issue tracking, documentation, source code organization, testing practices, dependencies, and community interaction. Public documentation and project websites can add evidence about installation procedures, usage patterns, supported platforms, and design goals.

Repository-based measurement has inherent limitations. Repository metadata represents a snapshot in time: last commit dates, issue counts, release information, and pull request activity can change after measurement. Some projects distribute their development activity across multiple repositories or use external issue trackers, making it difficult to capture a complete picture from any single repository. The interpretation of repository metrics requires care: a high number of open issues may indicate an active user community rather than poor maintenance, and a low commit frequency may reflect project maturity rather than abandonment.

The measurement record therefore documents snapshot dates, repository choices, and interpretation rules. Proprietary tools and tools with closed core components may be important in robotics practice but fall outside this procedure (Section~\ref{sec:software-categories}); Chapter~\ref{ch:methodology} discusses the specific inclusion and exclusion decisions that follow.
